\documentclass[main.tex]{subfiles}
\begin{document}

\chapter{The Likelihood Ratio Test and Wilks' Theorem}\label{ch:lrt}\chaptermark{The Likelihood Ratio Test}

UMP tests are elegant but rare. For most hypotheses met in practice --- composite null hypotheses, nuisance parameters, multidimensional constraints --- a general method is needed. The generalised likelihood ratio test (LRT) extends the Neyman--Pearson ratio naturally to composite hypotheses: it compares the best fits of the restricted and unrestricted models. Wilks' theorem says that, under $H_0$, this test statistic converges to a chi-square distribution independently of the model. Together with the Wald and score tests, the LRT forms the ``holy trinity'' of asymptotic tests; almost all tests in regression, GLMs and goodness of fit are special cases of this trinity.

\section{Definitions}

\begin{definition}[Generalised Likelihood Ratio]
For $H_0:\bm\theta\in\Theta_0$, $H_1:\bm\theta\in\Theta\setminus\Theta_0$,
\[
\Lambda(\x)=\frac{\sup_{\bm\theta\in\Theta_0}L(\bm\theta;\x)}{\sup_{\bm\theta\in\Theta}L(\bm\theta;\x)}=\frac{L(\hat{\bm\theta}_0;\x)}{L(\hat{\bm\theta};\x)}\in[0,1]
\]
is the \emph{generalised likelihood ratio} ($\hat{\bm\theta}_0$: restricted MLE, $\hat{\bm\theta}$: unrestricted MLE). The \emph{likelihood ratio test} (LRT) rejects $H_0$ when $\Lambda(\x)\le k$. Equivalently, it rejects when
\[
G^2=-2\ln\Lambda=2\big[\ell(\hat{\bm\theta})-\ell(\hat{\bm\theta}_0)\big]\ge0
\]
is large.
\end{definition}

\begin{definition}[Wald and Score (Rao) Tests]
For the simple $H_0:\theta=\theta_0$ (one parameter),
\[
W=\big(\hat\theta-\theta_0\big)^2\,I_n(\hat\theta)\qquad\text{(Wald)},\qquad
S=\frac{\big[\ell'(\theta_0)\big]^2}{I_n(\theta_0)}\qquad\text{(score / Lagrange multiplier)}.
\]
In the multidimensional case $W=(\hat{\bm\theta}-\bm\theta_0)^\top I_n(\hat{\bm\theta})(\hat{\bm\theta}-\bm\theta_0)$ and $S=\nabla\ell(\bm\theta_0)^\top I_n(\bm\theta_0)^{-1}\nabla\ell(\bm\theta_0)$. The Wald test requires fitting only the unrestricted model, the score test only the restricted one; the LRT requires both.
\end{definition}

\section{Main Theorems and Proofs}

\begin{proposition}[The LRT Is Based on a Sufficient Statistic]
If $T$ is sufficient, then $\Lambda(\x)$ is a function of $T(\x)$ alone.
\end{proposition}
\begin{proof}
By factorisation $L(\bm\theta;\x)=g(T(\x),\bm\theta)h(\x)$; $h(\x)$ cancels between numerator and denominator.
\end{proof}

\begin{proposition}[For Simple Hypotheses, LRT = NP Test]
If $\Theta=\{\theta_0,\theta_1\}$, then
\[
\Lambda=\min\Big\{1,\,\frac{L(\theta_0)}{L(\theta_1)}\Big\},
\]
and $\Lambda\le k<1$ and $L(\theta_1)/L(\theta_0)\ge1/k$ give the same rejection region. The LRT is the natural generalisation of the Neyman--Pearson test.
\end{proposition}
\begin{proof}
Since $\Theta_0=\{\theta_0\}$ and $\Theta=\{\theta_0,\theta_1\}$,
\[
\Lambda=\frac{L(\theta_0)}{\max\{L(\theta_0),L(\theta_1)\}}=\min\Big\{1,\frac{L(\theta_0)}{L(\theta_1)}\Big\}.
\]
 For $k<1$, $\Lambda\le k$ can occur only when the minimum comes from the second term, which means $L(\theta_0)/L(\theta_1)\le k$, i.e.\ $L(\theta_1)/L(\theta_0)\ge1/k$.
\end{proof}

\begin{theorem}[Wilks' Theorem \cite{wilks38}]\label{thm:wilks}
Let $X_i\iid f(\cdot;\bm\theta)$ with $\Theta\subseteq\R^k$ open, and suppose the conditions for asymptotic normality of the MLE of Chapter~\ref{ch:mle-asimptotik} hold. Let $\Theta_0$ be a smooth $(k-r)$-dimensional submanifold of $\Theta$ (e.g.\ a subspace in which $r$ components are fixed), and let the true $\bm\theta_0$ be an interior point of $\Theta_0$. Then under $H_0$
\[
G_n^2 = -2 \ln \Lambda_n(\X) \dto \chi^2_r,\qquad r = \dim(\Theta) - \dim(\Theta_0).
\]
\end{theorem}
\begin{proof}[Proof (simple $H_0$, $k=r=1$)]
Expand $\ell(\theta_0)$ in a Taylor series around $\hat\theta$; since $\ell'(\hat\theta)=0$, for some $\theta^*$ between the two,
\[
\ell(\theta_0)=\ell(\hat\theta)+\tfrac12(\theta_0-\hat\theta)^2\ell''(\theta^*).
\]
Hence
\[
G^2=2[\ell(\hat\theta)-\ell(\theta_0)]=\Big[\sqrt n(\hat\theta-\theta_0)\Big]^2\cdot\Big(-\frac1n\ell''(\theta^*)\Big).
\]
By Theorem~\ref{thm:mle-an}, $\sqrt n(\hat\theta-\theta_0)\sqrt{I_1(\theta_0)}\dto\Normal(0,1)$; since the MLE is consistent, $\theta^*\pto\theta_0$ and (by the third-derivative condition) $-\frac1n\ell''(\theta^*)\pto I_1(\theta_0)$. By Slutsky and the continuous mapping theorem, $G^2=\big[\sqrt{nI_1}(\hat\theta-\theta_0)\big]^2\cdot\frac{-\ell''(\theta^*)/n}{I_1}\dto Z^2\cdot1=\chi^2_1$.

In the general case the same expansion gives $G^2\approx\sqrt n(\hat{\bm\theta}-\hat{\bm\theta}_0)^\top I_1\sqrt n(\hat{\bm\theta}-\hat{\bm\theta}_0)$; this is the squared norm of the $r$-dimensional component of an $\Normal_k(\mathbf 0,I_1^{-1})$ vector orthogonal (in the Fisher metric) to the tangent space of $\Theta_0$, and it has a $\chi^2_r$ distribution.
\end{proof}

\begin{theorem}[Asymptotic Equivalence: LRT, Wald, Score]\label{thm:uclu}
Under the conditions of Wilks' theorem and when $H_0$ is true, $W_n-G_n^2\pto0$ and $S_n-G_n^2\pto0$; hence all three statistics have the limit distribution $\chi^2_r$. Moreover, under $\sqrt n$-local alternatives $\bm\theta_n=\bm\theta_0+\mathbf h/\sqrt n$ all three converge to the same non-central $\chi^2_r(\delta)$ distribution; that is, their asymptotic powers are equal.
\end{theorem}
\begin{proof}[Idea of proof ($r=k=1$)]
From the proof of Wilks' theorem, $G^2=n(\hat\theta-\theta_0)^2I_1+o_P(1)$. Since $W=n(\hat\theta-\theta_0)^2I_1(\hat\theta)$ and $I_1(\hat\theta)\pto I_1(\theta_0)$, $W-G^2\pto0$. For the score, $\ell'(\theta_0)=-\ell''(\theta^{**})(\hat\theta-\theta_0)\approx nI_1(\hat\theta-\theta_0)$, so $S=\ell'(\theta_0)^2/(nI_1)\approx n I_1(\hat\theta-\theta_0)^2$.
\end{proof}

\begin{remark}[The Assumptions of Wilks' Theorem]
\begin{itemize}
  \item \textbf{Interior point:} if $\bm\theta_0$ lies on the boundary of $\Theta$ (e.g.\ $H_0:\sigma^2_{\text{random effect}}=0$), the limit is the mixture $\frac12\chi^2_0+\frac12\chi^2_1$; using $\chi^2_1$ gives a conservative test.
  \item \textbf{Identifiability:} if some parameters become unidentified under $H_0$ (e.g.\ testing the number of components in a mixture model), Wilks' theorem fails.
  \item \textbf{Fixed dimension:} if the number of parameters grows with $n$, the chi-square approximation may break down.
\end{itemize}
\end{remark}

\section{Worked Examples}

\begin{example}[Poisson Mean]\label{ex:poisson-lrt}
$X_i\iid\text{Poisson}(\lambda)$, $H_0:\lambda=\lambda_0$. Since $\hat\lambda=\bar x$ and $\ell(\lambda)=n\bar x\ln\lambda-n\lambda+\text{const}$,
\[
G^2=2n\Big[\bar x\ln\frac{\bar x}{\lambda_0}-(\bar x-\lambda_0)\Big],\qquad
W=\frac{n(\bar x-\lambda_0)^2}{\bar x},\qquad S=\frac{n(\bar x-\lambda_0)^2}{\lambda_0}.
\]
(If $\bar x=0$, the convention $0\ln0=0$ gives $G^2=2n\lambda_0$; the Wald statistic is then undefined.) By Wilks, $G^2\dto\chi^2_1$. The score statistic is the counterpart of the one-cell Pearson chi-square statistic.
\end{example}

\begin{example}[Normal with Unknown Variance: LRT $=$ $t$-Test]\label{ex:lrt-t}
$X_i\iid\Normal(\mu,\sigma^2)$, $H_0:\mu=\mu_0$, with $\sigma^2$ a nuisance parameter. The unrestricted MLE is $(\bar x,\hat\sigma^2)$, the restricted MLE $(\mu_0,\hat\sigma_0^2)$ with $\hat\sigma_0^2=\frac1n\sum(x_i-\mu_0)^2=\hat\sigma^2+(\bar x-\mu_0)^2$. Since the exponential term at the maximum of both likelihoods is $e^{-n/2}$,
\[
\Lambda=\Big(\frac{\hat\sigma^2}{\hat\sigma_0^2}\Big)^{n/2}=\Big(1+\frac{(\bar x-\mu_0)^2}{\hat\sigma^2}\Big)^{-n/2}=\Big(1+\frac{t^2}{n-1}\Big)^{-n/2},\qquad t=\frac{\sqrt n(\bar x-\mu_0)}{S}.
\]
Since $\Lambda$ is a decreasing function of $|t|$, the LRT is the two-sided $t$-test, and its distribution under $H_0$ is known \emph{exactly} ($t_{n-1}$; Chapter~\ref{ch:guven}). As $n\to\infty$, $G^2=n\ln(1+t^2/(n-1))\approx t^2\dto\chi^2_1$, consistent with Wilks.
\end{example}

\begin{example}[Multinomial Goodness of Fit: $G^2$ and Pearson's $X^2$ \cite{pearson1900}]\label{ex:multinom}
Let the cell counts $(N_1,\dots,N_k)$ have a $\text{Multinom}(n;p_1,\dots,p_k)$ distribution and test $H_0:p_j=p_{j0}$ ($j=1,\dots,k$). The unrestricted MLE is $\hat p_j=N_j/n$; $\dim\Theta=k-1$, $\dim\Theta_0=0$. Thus
\[
G^2=2\sum_{j=1}^kN_j\ln\frac{N_j}{np_{j0}}\dto\chi^2_{k-1},\qquad
X^2=\sum_{j=1}^k\frac{(N_j-np_{j0})^2}{np_{j0}}\dto\chi^2_{k-1}.
\]
$X^2$ is the score test for this problem; using the expansion $\ln(1+u)\approx u-u^2/2$ one can show that $G^2=X^2+o_P(1)$. If $H_0$ involves estimating $s$ parameters, the degrees of freedom become $k-1-s$.
\end{example}

\begin{example}[Equality of Two Poisson Means]
$X_1,\dots,X_m\iid\text{Poisson}(\lambda_1)$, $Y_1,\dots,Y_n\iid\text{Poisson}(\lambda_2)$, $H_0:\lambda_1=\lambda_2$. $\dim\Theta=2$, $\dim\Theta_0=1$, $r=1$. The restricted MLE is $\hat\lambda=\frac{\sum X_i+\sum Y_j}{m+n}$ and
\[
G^2=2\Big[S_X\ln\frac{\bar x}{\hat\lambda}+S_Y\ln\frac{\bar y}{\hat\lambda}\Big]\dto\chi^2_1,\qquad S_X=\textstyle\sum X_i,\ S_Y=\sum Y_j.
\]
\end{example}

\section{Computational Applications}

\subsection{Verifying Wilks' Theorem (Poisson)}

\pycode{code/ch15/c15_1_verifying_wilks_theorem_poisson.py}
\begin{lstlisting}[style=output]
E[G^2]: 0.9998 (theory: 1.0) | type I error: 0.0517
Var[G^2]: 2.0141 (theory: 2.0) | KS p-value (against chi2_1): 0.000
\end{lstlisting}

The very small $p$-value of the KS test arises because, in an enormous simulation with $150\,000$ replications, it detects even the small deviations caused by the discreteness of $G^2$ (a Poisson sum is an integer); the mean, the variance and the type I error agree with $\chi^2_1$. This is a nice example of the difference between ``statistical significance'' and ``practical importance'' in large-sample simulations.

\subsection{The Holy Trinity: Exact Sizes in Small Samples}
Since the Poisson sum $T=n\bar X\sim\text{Poisson}(n\lambda_0)$, the true sizes of the three tests under $H_0$ can be computed exactly, without simulation.

\pycode{code/ch15/c15_2_holy_trinity_exact_sizes.py}
\begin{lstlisting}[style=output]
 n   lambda0 |  LRT    Wald   Score  (target 0.05)
   5    1.0   | 0.0723 0.1301 0.0386
  10    1.0   | 0.0563 0.0743 0.0374
  10    0.3   | 0.0833 0.2030 0.0335
  50    1.0   | 0.0474 0.0512 0.0556
 200    1.0   | 0.0476 0.0537 0.0517
\end{lstlisting}

For small values of $n\lambda_0$ the Wald test badly exceeds the size ($20\%$ instead of $5\%$ for $n\lambda_0=3$!): the probability of $\bar x=0$ is high, and in that case $W=\infty$. The LRT also exceeds the nominal level, but much more mildly; the score test is slightly conservative in this example. The practical lesson: in small samples the Wald test should be avoided and the score test or the LRT preferred. As $n$ grows, all three approach the nominal level (Theorem~\ref{thm:uclu}). Because of the discreteness, the sizes do not vary monotonically with $n$ but fluctuate in a saw-tooth pattern.

\subsection{Multinomial Goodness of Fit}

\pycode{code/ch15/c15_3_multinomial_goodness_of_fit.py}
\begin{lstlisting}[style=output]
n=40: P(G^2 > chi2_3,0.05) = 0.0600 | P(X^2 > chi2_3,0.05) = 0.0474
means: G^2=3.104, X^2=2.996 (theory 3) | Corr(G^2, X^2)=0.9666
Observed [ 7  5 13 15], expected [ 4.  8. 12. 16.]
Pearson X^2 = 3.5208, p = 0.3181 | G^2 = 3.2795, p = 0.3505
\end{lstlisting}

With $n=40$ and a smallest expected cell frequency of $4$, the size of Pearson's $X^2$ is closer to the nominal level than that of $G^2$; for small expected frequencies the chi-square approximation for $G^2$ improves more slowly. The rule of thumb ``expected frequency $\ge5$ in every cell'' is based on this observation.

\section{Comments and Notes}

\begin{remark}[Which Test?]
\begin{itemize}
  \item The \textbf{LRT} is invariant under reparametrisation (Zehna invariance) and usually gives the best chi-square approximation in small samples; with the Bartlett correction its error can be reduced to $O(n^{-2})$.
  \item The \textbf{Wald} test requires only the unrestricted model and is readily available in standard output (coefficient/SE); however, it is sensitive to the parametrisation (e.g.\ the result changes if $\ln\theta$ is tested instead of $\theta$) and behaves badly near the boundary (the Hauck--Donner effect).
  \item The \textbf{score} test requires only the model under $H_0$; this makes it ideal for model-extension tests (e.g.\ should an extra variable be added?) and diagnostic tests. Pearson's $X^2$, the log-rank test and the Cochran--Armitage trend test are score tests.
\end{itemize}
\end{remark}

\begin{remark}[Connection with Confidence Regions]
All three tests can be inverted into confidence regions (Chapter~\ref{ch:guven}, Theorem~\ref{thm:dualite}): inverting the LRT gives the \emph{profile likelihood interval} $\{\theta:2[\ell(\hat\theta)-\ell(\theta)]\le\chi^2_{1,\alpha}\}$, inverting Wald gives the symmetric interval $\hat\theta\pm z_{\alpha/2}\mathrm{SE}$, and inverting the score test gives (for the binomial) the Wilson interval.
\end{remark}

% --- CHAPTER SUMMARY (generated by tools/summaries.py) ---
\section{Chapter Summary}
\begin{summary}
\begin{itemize}
  \item The generalised likelihood ratio $\Lambda=\sup_{\Theta_0}L/\sup_\Theta L$ depends only on a sufficient statistic. For simple hypotheses it reduces to the Neyman--Pearson test.
  \item Wilks' theorem (Theorem~\ref{thm:wilks}): in regular models, under $H_0$, $-2\ln\Lambda\dto\chi^2_r$, where $r$ is the difference in the numbers of free parameters.
  \item The likelihood ratio, Wald and score tests are asymptotically equivalent (Theorem~\ref{thm:uclu}). In finite samples they may give different results; the Wald test is not invariant under reparametrisation.
  \item If the true parameter lies on the boundary or the model is not identifiable (e.g.\ the number of mixture components), the Wilks approximation breaks down.
  \item In the multinomial model the goodness-of-fit statistics $G^2$ and Pearson's $X^2$ are special cases of the likelihood ratio approach.
\end{itemize}
\end{summary}
% --- END OF CHAPTER SUMMARY ---

\paragraph{Further reading.} \cite[Section~10.3]{casella}, \cite[Chapters~16--17]{vdv}, \cite[Chapter~22]{ferguson}.

\section{Exercises}

\begin{exercise}\label{exr:15-1}
$X_i\iid\Normal(\mu,\sigma^2)$, $H_0:\sigma^2=\sigma_0^2$ ($\mu$ unknown). Write $G^2$ in terms of $S^2/\sigma_0^2$ and show that the LRT is a chi-square test that is \emph{not equal-tailed}.
\hfill\hyperref[coz:15-1]{\footnotesize$\rightarrow$\,Solution}
\end{exercise}

\begin{exercise}\label{exr:15-2}
Write down the $G^2$ and Pearson $X^2$ statistics for the hypothesis of independence in a $2\times2$ contingency table. Verify that the number of degrees of freedom is $(2-1)(2-1)=1$ by computing $\dim\Theta-\dim\Theta_0$.
\hfill\hyperref[coz:15-2]{\footnotesize$\rightarrow$\,Solution}
\end{exercise}

\begin{exercise}\label{exr:15-3}
For the exponential distribution and $H_0:\lambda=\lambda_0$, construct the Wald test (a) on the $\lambda$ scale and (b) on the $\ln\lambda$ scale. For $n=10$, compute and compare the exact sizes of the two tests (using $2\lambda_0\sum X_i\sim\chi^2_{2n}$). Compare them with the LRT.
\hfill\hyperref[coz:15-3]{\footnotesize$\rightarrow$\,Solution}
\end{exercise}

\begin{exercise}\label{exr:15-4}
A boundary parameter: $X_i\iid\Normal(\mu,1)$, $\Theta=[0,\infty)$, $H_0:\mu=0$. Show that $G^2=n\bar X^2\I\{\bar X>0\}$ and that under $H_0$, $G^2\dto\frac12\chi^2_0+\frac12\chi^2_1$. What is the true size of the test that uses the $\chi^2_1$ critical value?
\hfill\hyperref[coz:15-4]{\footnotesize$\rightarrow$\,Solution}
\end{exercise}

\end{document}
